This appendix collects the algebraic details behind
Section~\ref{sec:method}. These derivations are used to motivate testable
predictions, not to prove full Neural Collapse convergence for deep networks.

\subsection{Coupled and Decoupled Weight Decay}
\label{app:coupled-decoupled-derivation}

For a parameter block \(\theta\), let the coupled and decoupled components of a
fixed total weight decay be
\[
    \lambda_{\mathrm{coupled}}
    =
    \gamma\lambda_{\mathrm{total}},
    \qquad
    \lambda_{\mathrm{decoupled}}
    =
    (1-\gamma)\lambda_{\mathrm{total}}.
\]
The coupled component enters the gradient,
\[
    g_t^{(\gamma)}
    =
    \nabla_\theta \mathcal{L}(\theta_t)
    +
    \lambda_{\mathrm{coupled}}\theta_t,
\]
whereas the decoupled component is applied as direct parameter shrinkage:
\[
    \theta_{t+1}
    =
    (1-\eta_t\lambda_{\mathrm{decoupled}})\theta_t
    -
    \eta_t\mathcal{P}_t(g_t^{(\gamma)}).
\]
Here \(\mathcal{P}_t\) denotes the optimizer-specific momentum or preconditioning
operator.

For a non-adaptive optimizer without momentum, coupled L2 decay and decoupled
weight decay induce the same first-order update:
\[
    \theta_{t+1}
    =
    \theta_t-\eta_t\nabla_\theta\mathcal{L}(\theta_t)-\eta_t\lambda\theta_t.
\]
Under an adaptive diagonal preconditioner \(D_t\), however, Adam and AdamW differ
in whether the regularization term is preconditioned:
\[
    \theta_{t+1}^{\mathrm{Adam}}
    \approx
    \theta_t-\eta_tD_t\bigl(\nabla_\theta\mathcal{L}(\theta_t)+\lambda\theta_t\bigr),
\]
whereas
\[
    \theta_{t+1}^{\mathrm{AdamW}}
    \approx
    (1-\eta_t\lambda)\theta_t-\eta_tD_t\nabla_\theta\mathcal{L}(\theta_t).
\]
Thus Adam applies weight decay in the preconditioned geometry, while AdamW
applies weight decay in raw parameter space.

\subsection{Local Residual and Covariance Dynamics}
\label{app:local-residual-covariance-dynamics}

We connect the optimizer update to penultimate feature geometry through a local
residual approximation. Let
\[
    h_i = \mu_{y_i}+\varepsilon_i
\]
be the decomposition of a training feature into its class mean and residual.
Consider a terminal-phase or locally linearized regime in which the feature
residual of a class-\(c\) sample evolves approximately under a local curvature
operator \(B_{c,t}\):
\[
    \nabla_{h_i}\mathcal{L}
    \approx
    B_{c,t}\varepsilon_{i,t}.
\]
This approximation is exact in a simple peeled-feature quadratic model and is
used here only as a first-order description of detector-relevant residual
geometry.

Under this approximation, the residual update takes different forms under
different optimizers. For SGD or SGDW,
\[
    \varepsilon_{i,t+1}^{\mathrm{SGD/SGDW}}
    \approx
    \bigl(I-\eta_t(B_{c,t}+\lambda I)\bigr)\varepsilon_{i,t}.
\]
For coupled Adam,
\[
    \varepsilon_{i,t+1}^{\mathrm{Adam}}
    \approx
    \bigl(I-\eta_tD_t^h(B_{c,t}+\lambda I)\bigr)\varepsilon_{i,t},
\]
where \(D_t^h\) is the effective feature-side preconditioner induced by the
optimizer and the local Jacobian. For AdamW,
\[
    \varepsilon_{i,t+1}^{\mathrm{AdamW}}
    \approx
    \bigl(I-\eta_t(\lambda I+D_t^hB_{c,t})\bigr)\varepsilon_{i,t}.
\]

Let \(M_{c,t}\) denote the corresponding residual transition matrix. Then the
within-class covariance evolves approximately as
\[
    \Sigma_{c,t+1}
    \approx
    M_{c,t}\Sigma_{c,t}M_{c,t}^{\top}.
\]
The transition matrices are
\[
    M_{c,t}^{\mathrm{SGD/SGDW}}
    =
    I-\eta_t(B_{c,t}+\lambda I),
\]
\[
    M_{c,t}^{\mathrm{Adam}}
    =
    I-\eta_tD_t^h(B_{c,t}+\lambda I),
\]
and
\[
    M_{c,t}^{\mathrm{AdamW}}
    =
    I-\eta_t(\lambda I+D_t^hB_{c,t}).
\]
Consequently, Adam and AdamW can contract the same residual directions at
different rates. These direction-dependent contraction rates can change
within-class dispersion, covariance eigenspectra, anisotropy, effective rank,
and feature-norm dispersion, even when the final ID accuracy is similar.

\subsection{Covariance Expansion and Precision-Weighted Scores}
\label{app:covariance-expansion-precision}

Mahalanobis scoring and Gaussian feature-density scoring both contain
precision-weighted quadratic terms. For Mahalanobis, the relevant covariance is
typically a tied within-class covariance. For class-conditional Gaussian density,
the relevant covariance is the class covariance.

If \(A^{(1)}\succeq A^{(0)}\succ 0\), then matrix inversion reverses the
positive-semidefinite order:
\[
    (A^{(1)})^{-1}\preceq (A^{(0)})^{-1}.
\]
Therefore, for any residual vector \(\delta\),
\[
    \delta^\top(A^{(1)})^{-1}\delta
    \le
    \delta^\top(A^{(0)})^{-1}\delta.
\]
Thus, if an optimizer expands detector-relevant covariance directions, samples
aligned with those directions receive smaller Mahalanobis or Gaussian quadratic
penalties and can become more ID-like under feature-density scores.

The effect is direction dependent. In whitened coordinates, write
\[
    A^{(1)}
    =
    (A^{(0)})^{1/2}
    U\operatorname{diag}(\beta_i)U^\top
    (A^{(0)})^{1/2},
    \qquad
    \beta_i\ge 1.
\]
With \(\widetilde{\delta}=U^\top(A^{(0)})^{-1/2}\delta\),
\[
    \delta^\top(A^{(1)})^{-1}\delta
    =
    \sum_i
    \frac{\widetilde{\delta}_i^2}{\beta_i}.
\]
The reduction from the original penalty is
\[
    \sum_i
    \left(1-\frac{1}{\beta_i}\right)
    \widetilde{\delta}_i^2.
\]
Thus two samples can shift by different amounts depending on their alignment
with expanded covariance directions. This is why covariance changes can alter
OOD ranking rather than merely rescale all scores uniformly.
